% Tabla: Resumen de memoria Flash/RAM — v0.16.0 vs v0.17.0
% Fecha: 2026-03-04

\begin{table}[htbp]
\centering
\caption{Comparación de uso de memoria — v0.16.0 vs v0.17.0}
\label{tab:memory-v017}
\begin{tabular}{l r r r r}
\toprule
\textbf{Recurso} & \textbf{Disponible} & \textbf{v0.16.0} & \textbf{v0.17.0} & \textbf{$\Delta$} \\
\midrule
Flash (bytes) & 4\,194\,176 & 709\,636 (16.91\,\%) & 644\,200 (15.36\,\%) & $-$65\,436 ($-$1.55\,pp) \\
RAM   (bytes) & 488\,976   & 312\,068 (63.82\,\%) & 317\,248 (64.88\,\%) & $+$5\,180 ($+$1.06\,pp) \\
\bottomrule
\end{tabular}

\smallskip
\footnotesize
\textit{Nota:} La reducción de Flash se debe a la eliminación de la pre-copia de
\texttt{last\_good} en \texttt{meter\_read\_all()} y optimización de flujo de control.
El incremento de RAM corresponde al campo \texttt{field\_mask} (\texttt{uint32\_t}) y
\texttt{read\_target} (\texttt{int}) añadidos a \texttt{struct meter\_readings}, más
la variable \texttt{consecutive\_meter\_failures} en \texttt{main.c}.
\end{table}
